% GENERATED FILE. Edit canonical lessons, then run npm run generate:books.
% canonical-source-hash: fnv1a64:e59309d27008f93f
% canonical-lessons: ES-C643-derrapar, ES-C643-arrancar, ES-C643-atascarse, ES-C643-acelerar, ES-C643-circular

\chapter{Al volante}
\label{ch:es-al-volante}
\hlchaptermodality{643}

\begin{chapteropening}
\textbf{By the end of this chapter:} \emph{I can say what a car does: start, skid, get stuck, speed up and move with the traffic.}

Complete the last lesson of chapter 643: I can say what a car does: start, skid, get stuck, speed up and move with the traffic.
\end{chapteropening}

\section[derrapar]{derrapar — to skid}

\label{lesson:ES-C643-derrapar}



\begin{quote}
Before the new one: say the Spanish for the pothole, then the Spanish for the stretch.
\end{quote}

\subsection*{You'll want to know: derrapar}
\textbf{derrapar} — \textquotedblleft{}to skid\textquotedblright{}.

\begin{quote}
Con la lluvia, el coche derrapó.
\end{quote}

\begin{quote}
\emph{The car skidded in the rain.}
\end{quote}

\emph{Un derrape} is a skid.

\begin{cousinweb}
\emph{derrapar} was borrowed from French \emph{déraper} (EtymDB), first said of an anchor that comes loose from the sea bed.
\end{cousinweb}

\begin{grammarlens}[title={What you've built}]
One more word for telling, explaining and describing at B1.
\end{grammarlens}

\subsection*{Your turn}
\begin{itemize}
\raggedright
  \item \emph{Say it:} \emph{derrapar}
  \item \emph{Say it:} \emph{derrapar}, once more
  \item \emph{Say it:} a sentence of your own with \emph{derrapar}
  \item \emph{Recall:} say the Spanish for the pothole, then the Spanish for the stretch, then say \emph{derrapar} again
\end{itemize}

\begin{tcolorbox}[breakable,colback=teal!4,colframe=teal!35!black,title={Before you move on}]
What does derrapar mean? (\textquotedblleft{}To skid\textquotedblright{}.) Say it once more.
\end{tcolorbox}
\section[arrancar]{arrancar — to start (an engine); to pull out, to tear out}

\label{lesson:ES-C643-arrancar}



\begin{quote}
Before the new one: say the Spanish for the stretch, then the Spanish for to skid.
\end{quote}

\subsection*{You'll want to know: arrancar}
\textbf{arrancar} — \textquotedblleft{}to start (an engine); to pull out, to tear out\textquotedblright{}.

\begin{quote}
El coche no arranca.
\end{quote}

\begin{quote}
\emph{The car will not start.}
\end{quote}

The older sense is pulling out by force: \emph{arrancar una hoja} — to tear out a page.

\begin{cousinweb}
The origin of \emph{arrancar} is not settled. One proposal links it to Latin \emph{ēruncāre} \textquotedblleft{}to pull up weeds\textquotedblright{}, but the path is disputed.
\end{cousinweb}

\begin{grammarlens}[title={What you've built}]
One more word for telling, explaining and describing at B1.
\end{grammarlens}

\subsection*{Your turn}
\begin{itemize}
\raggedright
  \item \emph{Say it:} \emph{arrancar}
  \item \emph{Say it:} \emph{arrancar}, once more
  \item \emph{Say it:} a sentence of your own with \emph{arrancar}
  \item \emph{Recall:} say the Spanish for the stretch, then the Spanish for to skid, then say \emph{arrancar} again
\end{itemize}

\begin{tcolorbox}[breakable,colback=teal!4,colframe=teal!35!black,title={Before you move on}]
What does arrancar mean? (\textquotedblleft{}To start (an engine); to pull out, to tear out\textquotedblright{}.) Say it once more.
\end{tcolorbox}
\section[atascarse]{atascarse — to get stuck; to get jammed}

\label{lesson:ES-C643-atascarse}



\begin{quote}
Before the new one: say the Spanish for to skid, then the Spanish for to start.
\end{quote}

\subsection*{You'll want to know: atascarse}
\textbf{atascarse} — \textquotedblleft{}to get stuck; to get jammed\textquotedblright{}.

\begin{quote}
El coche se atascó en la arena.
\end{quote}

\begin{quote}
\emph{The car got stuck in the sand.}
\end{quote}

\emph{Un atasco} is a traffic jam: the same family.

\begin{cousinweb}
The origin of \emph{atascarse} is not settled; it is the verb behind \emph{atasco} \textquotedblleft{}traffic jam\textquotedblright{}.
\end{cousinweb}

\begin{grammarlens}[title={What you've built}]
One more word for telling, explaining and describing at B1.
\end{grammarlens}

\subsection*{Your turn}
\begin{itemize}
\raggedright
  \item \emph{Say it:} \emph{atascarse}
  \item \emph{Say it:} \emph{atascarse}, once more
  \item \emph{Say it:} a sentence of your own with \emph{atascarse}
  \item \emph{Recall:} say the Spanish for to skid, then the Spanish for to start, then say \emph{atascarse} again
\end{itemize}

\begin{tcolorbox}[breakable,colback=teal!4,colframe=teal!35!black,title={Before you move on}]
What does atascarse mean? (\textquotedblleft{}To get stuck; to get jammed\textquotedblright{}.) Say it once more.
\end{tcolorbox}
\section[acelerar]{acelerar — to speed up, to accelerate}

\label{lesson:ES-C643-acelerar}



\begin{quote}
Before the new one: say the Spanish for to start, then the Spanish for to get stuck.
\end{quote}

\subsection*{You'll want to know: acelerar}
\textbf{acelerar} — \textquotedblleft{}to speed up, to accelerate\textquotedblright{}.

\begin{quote}
No aceleres en la ciudad.
\end{quote}

\begin{quote}
\emph{Do not speed up in town.}
\end{quote}

The pedal is \emph{el acelerador}.

\begin{cousinweb}
\emph{acelerar} is a learned borrowing from Latin \emph{accelerāre} \textquotedblleft{}to hasten\textquotedblright{} (Lewis \& Short; EtymDB), built from \emph{ad} and \emph{celer} \textquotedblleft{}swift\textquotedblright{}.
\end{cousinweb}

\begin{grammarlens}[title={What you've built}]
One more word for telling, explaining and describing at B1.
\end{grammarlens}

\subsection*{Your turn}
\begin{itemize}
\raggedright
  \item \emph{Say it:} \emph{acelerar}
  \item \emph{Say it:} \emph{acelerar}, once more
  \item \emph{Say it:} a sentence of your own with \emph{acelerar}
  \item \emph{Recall:} say the Spanish for to start, then the Spanish for to get stuck, then say \emph{acelerar} again
\end{itemize}

\begin{tcolorbox}[breakable,colback=teal!4,colframe=teal!35!black,title={Before you move on}]
What does acelerar mean? (\textquotedblleft{}To speed up, to accelerate\textquotedblright{}.) Say it once more.
\end{tcolorbox}
\section[circular]{circular — to circulate, to drive (of traffic)}

\label{lesson:ES-C643-circular}



\begin{quote}
Before the new one: say the Spanish for to get stuck, then the Spanish for to speed up.
\end{quote}

\subsection*{You'll want to know: circular}
\textbf{circular} — \textquotedblleft{}to circulate, to drive (of traffic)\textquotedblright{}.

\begin{quote}
En España se circula por la derecha.
\end{quote}

\begin{quote}
\emph{In Spain you drive on the right.}
\end{quote}

A word for signs and rules: \emph{circule despacio} — drive slowly. As an adjective it means round.

\begin{cousinweb}
\emph{circular} is a learned borrowing from Late Latin \emph{circulāre} \textquotedblleft{}to make round\textquotedblright{} (EtymDB; Lewis \& Short), built on \emph{circulus} \textquotedblleft{}a small ring\textquotedblright{}, from \emph{circus} \textquotedblleft{}a ring\textquotedblright{}.
\end{cousinweb}

\begin{grammarlens}[title={What you've built}]
One more word for telling, explaining and describing at B1.
\end{grammarlens}

\subsection*{Your turn}
\begin{itemize}
\raggedright
  \item \emph{Say it:} \emph{circular}
  \item \emph{Say it:} \emph{circular}, once more
  \item \emph{Say it:} a sentence of your own with \emph{circular}
  \item \emph{Recall:} say the Spanish for to get stuck, then the Spanish for to speed up, then say \emph{circular} again
\end{itemize}

\begin{tcolorbox}[breakable,colback=teal!4,colframe=teal!35!black,title={Before you move on}]
What does circular mean? (\textquotedblleft{}To circulate, to drive (of traffic)\textquotedblright{}.) Say it once more.
\end{tcolorbox}
